\begin{theorem}[window-\texorpdfstring{$6$}{6} 输出层的 Ramanujan--Mellin 双对偶]\label{thm:conclusion-window6-output-ramanujan-mellin-duality}
沿用定理 \ref{thm:conclusion-window6-visible-crt-arithmetic-phase-space} 与定理 \ref{thm:conclusion-window6-hidden-multiplicity-neutrality-on-nonzero-strata} 的记号，记
\[
R_6:=\ZZ/(21\ZZ)\cong \{0,1,\dots,20\},
\qquad
d(r):=
\begin{cases}
4,& 0\le r\le 8,\\
3,& 9\le r\le 12,\\
2,& 13\le r\le 20,
\end{cases}
\]
并对每个 \(q\mid 21\) 定义
\[
M_q(s):=\sum_{r\in R_6}d(r)^s\,c_q(r),
\qquad
\mathcal C_q(T):=\sum_{r\in R_6}c_q(r)\min\{d(r),T\}.
\]
则对任意 \(q\mid 21\) 及任意 \(s\in\CC\) 满足 \(0<\Re s<1\)，有
\[
\boxed{\;
M_q(s)=s(1-s)\int_0^\infty T^{s-2}\mathcal C_q(T)\,\dd T.\;}
\]
反之，对任意 \(c\in(0,1)\) 与任意 \(T>0\)，有 Mellin--Barnes 反演公式
\[
\boxed{\;
\mathcal C_q(T)=\frac1{2\pi i}
\int_{c-i\infty}^{c+i\infty}
\frac{M_q(s)}{s(1-s)}\,T^{1-s}\,\dd s.\;}
\]
\end{theorem}
\begin{proof}
对任意整数 \(d\ge 1\) 与 \(0<\Re s<1\)，直接分段积分得
\[
\int_0^\infty T^{s-2}\min\{d,T\}\,\dd T
=
\int_0^d T^{s-1}\,\dd T+d\int_d^\infty T^{s-2}\,\dd T
=
\frac{d^s}{s(1-s)}.
\]
故
\[
d^s=s(1-s)\int_0^\infty T^{s-2}\min\{d,T\}\,\dd T.
\]
将该式乘以 \(c_q(r)\) 并对 \(r\in R_6\) 求和，即得
\[
M_q(s)
=
s(1-s)\int_0^\infty
T^{s-2}
\sum_{r\in R_6}c_q(r)\min\{d(r),T\}\,\dd T
=
s(1-s)\int_0^\infty T^{s-2}\mathcal C_q(T)\,\dd T.
\]

对反演公式，只需注意 \(M_q(s)\) 是有限个 \(d(r)^s\) 的线性组合，而 \(\mathcal C_q(T)\) 是对应的 \(\min\{d(r),T\}\) 的同一线性组合。主稿关于连续容量曲线与全阶矩谱的 Mellin--Barnes 完备互逆可逐项线性化应用，遂得所示反演式。证毕。
\end{proof}

\begin{theorem}[window-\texorpdfstring{$6$}{6} 三层多重度的 Ramanujan 电荷闭式]\label{thm:conclusion-window6-output-ramanujan-layer-charge-closedforms}
\leanpartial{paper\_conclusion\_window6\_output\_ramanujan\_layer\_charge\_closedforms}{The current Lean definition \texttt{window6RamanujanTerm} encodes hidden-template offsets rather than the residue Ramanujan sum \(c_q(r)\), so the requested signature does not match the checked code.}
记
\[
A_2:=\{13,\dots,20\},
\qquad
A_3:=\{9,10,11,12\},
\qquad
A_4:=\{0,\dots,8\}.
\]
则对导子 \(3,7,21\) 分别有层电荷向量
\[
\left(\sum_{r\in A_2}c_3(r),\ \sum_{r\in A_3}c_3(r),\ \sum_{r\in A_4}c_3(r)\right)=(-2,2,0),
\]
\[
\left(\sum_{r\in A_2}c_7(r),\ \sum_{r\in A_3}c_7(r),\ \sum_{r\in A_4}c_7(r)\right)=(-1,-4,5),
\]
\[
\left(\sum_{r\in A_2}c_{21}(r),\ \sum_{r\in A_3}c_{21}(r),\ \sum_{r\in A_4}c_{21}(r)\right)=(-5,-2,7).
\]
因而对一切 \(s\in\CC\)，成立闭式
\[
M_1(s)=8\cdot 2^s+4\cdot 3^s+9\cdot 4^s,
\]
\[
M_3(s)=2(3^s-2^s),
\qquad
M_7(s)=5\cdot 4^s-4\cdot 3^s-2^s,
\qquad
M_{21}(s)=7\cdot 4^s-2\cdot 3^s-5\cdot 2^s.
\]
\end{theorem}
\begin{proof}
由定理 \ref{thm:conclusion-window6-hidden-multiplicity-neutrality-on-nonzero-strata}，window-\(6\) 的纤维大小在 \(R_6\) 上恰分成三层
\[
d(r)=4\ \text{于 }A_4,\qquad d(r)=3\ \text{于 }A_3,\qquad d(r)=2\ \text{于 }A_2.
\]
同时由推论 \ref{cor:fold-bin-m6-fiber-hist}，
\[
M_1(s)=\sum_{r\in R_6}d(r)^s=8\cdot 2^s+4\cdot 3^s+9\cdot 4^s.
\]

对 \(q=3\)，利用
\[
c_3(r)=
\begin{cases}
2,& 3\mid r,\\
-1,& 3\nmid r,
\end{cases}
\]
逐层计数可得
\[
\sum_{r\in A_2}c_3(r)=-2,\qquad
\sum_{r\in A_3}c_3(r)=2,\qquad
\sum_{r\in A_4}c_3(r)=0.
\]
同理，对 \(q=7\)，由
\[
c_7(r)=
\begin{cases}
6,& 7\mid r,\\
-1,& 7\nmid r,
\end{cases}
\]
逐层计数即得
\[
\sum_{r\in A_2}c_7(r)=-1,\qquad
\sum_{r\in A_3}c_7(r)=-4,\qquad
\sum_{r\in A_4}c_7(r)=5.
\]

又因 \(3,7\) 互素且 Ramanujan 和满足乘法性，
\[
c_{21}(r)=c_3(r)c_7(r),
\]
故
\[
c_{21}(r)=
\begin{cases}
12,& 21\mid r,\\
-2,& 3\mid r,\ 7\nmid r,\\
-6,& 7\mid r,\ 3\nmid r,\\
1,& (r,21)=1.
\end{cases}
\]
逐层计数便得到 \((-5,-2,7)\)。

最后，把各导子的层电荷与三层多重度 \(2^s,3^s,4^s\) 配对求和，即得全部 \(M_q(s)\) 的闭式。证毕。
\end{proof}

\begin{theorem}[window-\texorpdfstring{$6$}{6} 截断容量轴上的导子激活律]\label{thm:conclusion-window6-output-truncated-capacity-conductor-activation}
对 \(T\ge 0\)，非平凡导子的截断容量变换满足分段闭式
\[
\mathcal C_3(T)=
\begin{cases}
0,& 0\le T\le 2,\\
2T-4,& 2\le T\le 3,\\
2,& T\ge 3,
\end{cases}
\]
\[
\mathcal C_7(T)=
\begin{cases}
0,& 0\le T\le 2,\\
T-2,& 2\le T\le 3,\\
5T-14,& 3\le T\le 4,\\
6,& T\ge 4,
\end{cases}
\]
\[
\mathcal C_{21}(T)=
\begin{cases}
0,& 0\le T\le 2,\\
5T-10,& 2\le T\le 3,\\
7T-16,& 3\le T\le 4,\\
12,& T\ge 4.
\end{cases}
\]
\end{theorem}
\begin{proof}
由定理 \ref{thm:conclusion-window6-output-ramanujan-layer-charge-closedforms}，
\[
\mathcal C_q(T)=\sum_{j=2}^4 L_{j,q}\min\{j,T\},
\]
其中 \((L_{2,q},L_{3,q},L_{4,q})\) 分别是导子 \(q\) 在三层 \((A_2,A_3,A_4)\) 上的电荷向量。于是
\[
\mathcal C_3(T)=-2\min\{2,T\}+2\min\{3,T\},
\]
\[
\mathcal C_7(T)=-\min\{2,T\}-4\min\{3,T\}+5\min\{4,T\},
\]
\[
\mathcal C_{21}(T)=-5\min\{2,T\}-2\min\{3,T\}+7\min\{4,T\}.
\]
在区间 \([0,2]\)、\([2,3]\)、\([3,4]\) 与 \([4,\infty)\) 上逐段化简，即得三组分段公式。证毕。
\end{proof}

\begin{corollary}[导子 \texorpdfstring{$3$}{3} 的中层纯探测性]\label{cor:conclusion-window6-output-conductor-three-middle-layer-probe}
记
\[
\mathcal N_q(T):=\frac{\dd}{\dd T}\mathcal C_q(T)
\qquad
\bigl(T\notin\{2,3,4\}\bigr).
\]
则
\[
\mathcal N_3(T)=
\begin{cases}
0,& 0<T<2,\\
2,& 2<T<3,\\
0,& 3<T,
\end{cases}
\]
\[
\mathcal N_7(T)=
\begin{cases}
0,& 0<T<2,\\
1,& 2<T<3,\\
5,& 3<T<4,\\
0,& 4<T,
\end{cases}
\]
\[
\mathcal N_{21}(T)=
\begin{cases}
0,& 0<T<2,\\
5,& 2<T<3,\\
7,& 3<T<4,\\
0,& 4<T.
\end{cases}
\]
因此，导子 \(3\) 在 window-\(6\) 的多重度压缩模型中是一个纯中层探针：它只读取 \(d=2\) 与 \(d=3\) 之间的过渡，而对顶层 \(d=4\) 严格失明。与之相反，导子 \(7\) 与 \(21\) 都能读取顶层，并且 \(21\) 的顶层权重严格大于 \(7\)。
\end{corollary}
\begin{proof}
对定理 \ref{thm:conclusion-window6-output-truncated-capacity-conductor-activation} 中的分段一次函数逐段求导即可。证毕。
\end{proof}

\begin{theorem}[window-\texorpdfstring{$6$}{6} 中合成导子 \texorpdfstring{$21$}{21} 的严格塌缩律]\label{thm:conclusion-window6-output-composite-conductor-collapse}
对一切 \(s\in\CC\) 与一切 \(T\ge 0\)，成立恒等式
\[
\boxed{\;
5M_{21}(s)=9M_3(s)+7M_7(s),\qquad
5\mathcal C_{21}(T)=9\mathcal C_3(T)+7\mathcal C_7(T).\;}
\]
等价地，在 window-\(6\) 的三层多重度压缩之后，合成导子 \(21\) 不再提供任何独立通道；它只是素导子 \(3\) 与 \(7\) 的严格线性影子。
\end{theorem}
\begin{proof}
由定理 \ref{thm:conclusion-window6-output-ramanujan-layer-charge-closedforms}，三层电荷向量满足逐层恒等式
\[
5(-5,-2,7)=9(-2,2,0)+7(-1,-4,5).
\]
把该恒等式作用到层基
\[
(2^s,3^s,4^s)
\]
上，即得
\[
5M_{21}(s)=9M_3(s)+7M_7(s).
\]
同理，把同一逐层恒等式作用到
\[
\bigl(\min\{2,T\},\min\{3,T\},\min\{4,T\}\bigr)
\]
上，即得
\[
5\mathcal C_{21}(T)=9\mathcal C_3(T)+7\mathcal C_7(T).
\]
证毕。
\end{proof}

\begin{theorem}[导子 \texorpdfstring{$(1,3,7)$}{(1,3,7)} 对三层多重度的完备可逆性]\label{thm:conclusion-window6-output-prime-conductor-complete-reconstruction}
记
\[
X(s):=
\begin{pmatrix}
2^s\\[2pt]
3^s\\[2pt]
4^s
\end{pmatrix},
\qquad
Y(s):=
\begin{pmatrix}
M_1(s)\\[2pt]
M_3(s)\\[2pt]
M_7(s)
\end{pmatrix}.
\]
则
\[
Y(s)=
\begin{pmatrix}
8&4&9\\
-2&2&0\\
-1&-4&5
\end{pmatrix}
X(s),
\qquad
\det
\begin{pmatrix}
8&4&9\\
-2&2&0\\
-1&-4&5
\end{pmatrix}
=210\neq 0.
\]
因此 \(X(s)\) 由 \((M_1(s),M_3(s),M_7(s))\) 唯一恢复；显式地，
\[
2^s=\frac{5M_1(s)-28M_3(s)-9M_7(s)}{105},
\]
\[
3^s=\frac{10M_1(s)+49M_3(s)-18M_7(s)}{210},
\]
\[
4^s=\frac{5M_1(s)+14M_3(s)+12M_7(s)}{105}.
\]
故在 window-\(6\) 的三层压缩模型上，平凡导子与两个素导子 \(3,7\) 已构成一组完备的 Ramanujan 层析坐标；任何合成导子信息都可由它们重建。
\end{theorem}
\begin{proof}
矩阵表达式正是定理 \ref{thm:conclusion-window6-output-ramanujan-layer-charge-closedforms} 中四个闭式的前三个所给出的线性关系。其行列式直接计算得
\[
210\neq 0,
\]
故该矩阵可逆。对 \(Y(s)\) 施加逆矩阵，即得所列三条显式反演公式。证毕。
\end{proof}

\begin{conjecture}[平方自由可见模上的素导子张成原理]\label{conj:conclusion-window6-output-squarefree-prime-conductor-span}
设某一分辨率层的可见模 \(N\) 为平方自由整数，并设对应多重度函数仅取有限多个值
\[
d_1,\dots,d_k.
\]
对每个 \(q\mid N\) 定义压缩 Ramanujan--moment
\[
M_q(s):=\sum_{r\in \ZZ/N\ZZ}d(r)^s\,c_q(r).
\]
则在由层函数 \((d_1^s,\dots,d_k^s)\) 张成的 \(k\)-维压缩空间中，应有
\[
\operatorname{span}\{M_q(s):q\mid N\}
=
\operatorname{span}\Bigl(\,M_1(s),\ \{M_p(s):p\mid N,\ p\text{ 为素数}\}\Bigr).
\]
换言之，一旦把微观几何压缩到多重度层，所有合成导子都应塌缩为素导子通道的线性阴影。window-\(6\) 的
\[
N=21=3\cdot 7
\]
与定理 \ref{thm:conclusion-window6-output-composite-conductor-collapse} 中的严格关系
\[
5M_{21}=9M_3+7M_7
\]
正是这一原则的最小非平凡样本。
\end{conjecture}
